\section{Further identities at degree \(163\)}
\label{sec:163-additional-identities}

The arithmetic square identity of
Lemma~\ref{lem:163-period-unit-square} has three useful refinements.
The first identifies precisely the index-four enlargement of the
circular-unit lattice needed to contain its square roots. The second
shows that a mixed beta ratio retains the same residual scalar
coordinate. The third places the square roots in an explicit
three-point orbit of a rational map. These identities do not select
an alternative in Proposition~\ref{prop:cm163-full-image-dichotomy}.

\subsection{The circular-unit saturation defect}

Put \(p=163\), \(\zeta=\zeta_p\), \(F^+=\mathbf Q(\zeta+\zeta^{-1})\),
and let \(\mathcal E\) be its real unit group.  Let \(C\subset\mathcal E\)
be the real circular-unit group, including \(-1\).  For nonzero
\(a\in\mathbf F_p\), use the standard real circular unit
\[
 \eta_a=\zeta^{(1-a)/2}\frac{1-\zeta^a}{1-\zeta},
\]
where the half exponent is interpreted modulo \(p\).
Then \(\eta_1=1\) and \(\eta_{-a}=-\eta_a\).  The group \(C\)
is generated by \(-1\) and the \(\eta_a\), taking one representative
from each pair \(\{a,-a\}\).
The circular-unit finite-index theorem gives
\(\operatorname{rank}(C/\{\pm1\})=80\); see
\cite[p.~206, the definition and index of \(C\)]{Dummit2020}.

Write
\[
 \epsilon_k=\zeta^{2^k}+\zeta^{-2^k}\qquad(0\leq k<81),
 \qquad
 U_i=\prod_{j=0}^{26}\epsilon_{i+3j}\quad(i=0,1,2).
\]
The equality \(2^{81}=-1\pmod {163}\) makes the \(\epsilon_k\)
periodic with period \(81\).  Let \(s(\zeta)=\zeta^2\).

\begin{proposition}[Circular-unit saturation defect]
\label{prop:163-circular-unit-saturation}
There is an isomorphism
\[
 C/\{\pm1\}\simeq
 \mathbf Z^{81}/\mathbf Z(1,\ldots,1)
\]
which sends the \(k\)-th standard vector to \(\epsilon_k\).
The classes of \(U_0,U_1,U_2\) in \(C/(\{\pm1\}C^2)\)
span the irreducible two-dimensional \(\mathbf F_2[C_3]\)-module;
their sole linear relation is
\[
 [U_0]+[U_1]+[U_2]=0.
\]
In particular none of the \(U_i\) is a square modulo sign of a
circular unit.

Let \(v_i=s^i(5+\theta_1)\), with the period notation of Lemma~\ref{lem:163-period-unit-square}.
Then \(U_i=-v_i^2\), and
\[
 C^*=\langle C,v_0,v_1,v_2\rangle\subset\mathcal E,
 \qquad C^*/C\simeq\mathbf F_2^2.
\]
The quotient has the same nontrivial irreducible \(C_3\)-action.
No assertion that \(C^*=\mathcal E\) is made.
\end{proposition}

\begin{proof}
The elementary identity
\[
 \frac{\eta_{4b}}{\eta_{2b}}
 =\zeta^{-b}\frac{1-\zeta^{4b}}{1-\zeta^{2b}}
 =\zeta^b+\zeta^{-b}
\]
gives
\(\epsilon_k=\eta_{2^{k+2}}/\eta_{2^{k+1}}\).
Modulo signs, the \(81\) classes \(2^k\) enumerate all nonzero
residues modulo sign.  Their consecutive differences therefore
generate all \(\eta_a\) modulo signs.  The \(80\) units
\(\eta_{2^k}\), \(1\leq k\leq80\), generate the free group
\(C/\{\pm1\}\) of rank \(80\), so they form a basis.
It follows that the only integral relation among the \(81\)
consecutive ratios, modulo signs, is their total product relation.
In fact telescoping gives the exact value
\[
 \prod_{k=0}^{80}\epsilon_k
 =\frac{\eta_{2^{82}}}{\eta_2}=-1.
\]
This proves the displayed presentation.

Modulo two, let \(w_i\in\mathbf F_2^{81}\) be the characteristic
vector of the indices congruent to \(i\pmod3\).  The image of
\(w_i\) is \([U_i]\).  The vector
\(a_0w_0+a_1w_1+a_2w_2\) is constant precisely when
\(a_0=a_1=a_2\), since the supports of the \(w_i\) are disjoint
and cover all \(81\) indices.  Their image therefore has dimension
two with exactly the claimed relation.  The automorphism \(s\)
cycles the three nonzero classes.  This is the irreducible nontrivial
\(C_3\)-module over \(\mathbf F_2\).

Lemma~\ref{lem:163-period-unit-square} gives \(v_i^2=-U_i\in C\).
Suppose \(v_0^{a_0}v_1^{a_1}v_2^{a_2}\in C\), with
\(a_i\in\{0,1\}\).  Squaring shows
\(U_0^{a_0}U_1^{a_1}U_2^{a_2}\in\{\pm1\}C^2\), so the
previous paragraph forces all \(a_i\) equal.  Conversely
\[
 (v_0v_1v_2)^2=(-1)^3U_0U_1U_2=1,
\]
and hence \(v_0v_1v_2=\pm1\in C\).  Thus the quotient is exactly
the stated group of order four.  The action is again the cyclic
permutation of its three nonzero directions.
\end{proof}

\paragraph{A concrete character separating the two notions of square.}
The preceding presentation permits a homomorphism
\(\ell:C\to\mathbf F_2\) with
\[
 \ell(-1)=0,\quad \ell(\epsilon_0)=\ell(\epsilon_1)=1,
 \quad\ell(\epsilon_k)=0\quad(2\leq k<81).
\]
The sum of these values is zero, as required by the sole relation
modulo signs.  This is a homomorphism on the circular-unit
group, so it respects every multiplicative identity among its
elements.  Nevertheless
\[
 (\ell(U_0),\ell(U_1),\ell(U_2))=(1,1,0).
\]
It cannot extend to a homomorphism \(C^*\to\mathbf F_2\):
an extension would send \(-v_0^2\) to zero.

Consequently an argument which only treats the \(\epsilon_k\)
as generators of the circular-unit group, respects their
multiplicative identities, and kills signs and circular squares
does not force the period products to vanish.  It must establish
the additional extension across the explicit index-four lattice
\(C\subset C^*\), or use other full-GT information.
This conclusion concerns the specified conductor and specified
multiplicative data.  It is not a construction of a GT cocycle,
not a counterexample to the full pentagon, and not a claim that
all higher-conductor distribution relations preserve this character.

\paragraph{Unit signatures.}
Dummit's Proposition~2 identifies a defect of order four between
totally positive circular units and circular squares at conductor
\(163\), while the full real unit group realizes all signatures
\cite[Proposition~2, pp.~209--210]{Dummit2020}.
The explicit lattice argument above isolates the three period products
and their square roots without importing a value for the full real
class number.  The equality of that class number to four is conditional on GRH
in the cited proposition and is not used here.

\subsection{A mixed beta ratio with the same residual coordinate}

Use the full coefficient-two beta values of a multiplier-one GT element
\(r\), and set \(p=163\), \(O=\mathbf Z_2[\zeta_p]\),
\(s(\zeta_p)=\zeta_p^2\).  For a unit of \(O\), write its
principal factor with a superscript \(+\); the complementary
Teichmuller factor has odd order.  All products below are in the
commutative unit group.

\begin{lemma}[Canonical principal gamma factor on an odd cyclic group]
\label{lem:finite-cyclic-principal-gamma}
Let \(p\) be odd and let
\(b:\mathbf F_p\times\mathbf F_p\to1+2O\) be a normalized
two-cocycle for the additive group \(\mathbf F_p\), with trivial
action on the coefficient group.  There is a unique function
\(\gamma:\mathbf F_p\to1+2O\), with \(\gamma(0)=1\), such that
\[
 b(a,b)=\frac{\gamma(a)\gamma(b)}{\gamma(a+b)}.
\]
Explicitly,
\[
 \gamma(a)=\left(\prod_{c\in\mathbf F_p}b(a,c)\right)^{1/p},
\]
where the \(p\)-th root is the unique root in the pro-two group
\(1+2O\).
\end{lemma}

\begin{proof}
Write \(T(a)=\prod_c b(a,c)\).  Multiplying the cocycle identity
\(b(a,b)b(a+b,c)=b(b,c)b(a,b+c)\) over \(c\) gives
\[
 b(a,b)^p T(a+b)=T(a)T(b).
\]
The \(p\)-power map is an automorphism of the abelian pro-two
group \(1+2O\), including its signs.  Taking its inverse proves
the claimed formula.  The ratio of two solutions would be a
homomorphism from a group of odd order to a pro-two group, hence
would be trivial.
\end{proof}

Apply the lemma to
\[
 b_r(a,b)=B_r(\zeta_p^a,\zeta_p^b)^+.
\]
The full-GT beta cocycle gives its hypothesis
\cite[Proposition~1.6.1]{Ihara1999},
\cite[Section~2.3, p.~101]{Furusho2003}. Frobenius
equivariance and beta unitarity give, by uniqueness,
\[
 \gamma_r(2a)=s(\gamma_r(a)),\qquad
 \gamma_r(-a)=\gamma_r(a)^{-1}.
\]
Thus \(u_r=\gamma_r(1)\) lies in the principal norm-one torus
\(T\).  If
\[
 a_r=B_r(\zeta_p,\zeta_p),\quad
 b_r=B_r(\zeta_p^2,\zeta_p),\quad
 k_r=a_r/b_r,
\]
then, writing group-ring actions additively,
\begin{equation}
 a_r^+=(2-s)u_r,
 \qquad k_r^+=(1+[3]-2s)u_r.                 \label{eq:163-gamma-beta-ratios}
\end{equation}

\begin{proposition}[Equality of the residual coordinates]
\label{prop:163-mixed-beta-residual}
For \(r\in R_{\rm global}\), the principal ratio \(k_r^+\)
belongs to \(P=P_{163}\), and
\[
 k_r^+\bmod QP=h_{163}(r)\bmod QP.
\]
Consequently the mixed ratio retains exactly the original
degree--163 alternative.  With \(n=2^{81}+1\) and the
first-layer detector \(\Psi:P\to W\) of
\eqref{eq:163-trace-detector},
\begin{equation}
 \Psi(k_r^n)=\Psi(h_{163}(r)).               \label{eq:163-mixed-ratio-detector}
\end{equation}
The expression on the left is defined because \(k_r^n\) is
principal and belongs to \(P\).
\end{proposition}

\begin{proof}
Equation~\eqref{eq:163-alexander-scalar} gives \(a_r^+=h_{163}(r)\in P\).
The operator \(2-s\) is an automorphism of \(T\), with inverse
\[
 (2-s)^{-1}=-\sum_{m\geq0}2^m s^{-m-1}.
\]
The series converges in the endomorphisms of the compact pro-two
group.  It preserves the closed stable subgroup \(P\).  Therefore
\(u_r\in P\), and \eqref{eq:163-gamma-beta-ratios} also places \(k_r^+\) in \(P\).

Direct modular arithmetic gives \(2^{20}=-3\pmod {163}\) and
\(2^{81}=-1\pmod {163}\), so
\[
 [3]=s^{101}.
\]
On \(E=P/QP\), the relations are \(2=0\), \(s^3=1\), and
\(s^2+s+1=0\).  Consequently
\[
 1+s^{101}-2s=1+s^2=s=2-s
                 \quad\text{on }E.
\]
Equation \eqref{eq:163-gamma-beta-ratios} proves the required equality in \(E\).

The full beta values in the ratio have norm one, by coefficient
equivariance and beta unitarity.  Hence the odd Teichmuller factor
of \(k_r\) has order dividing \(2^{81}+1=n\).  Thus
\(k_r^n=(k_r^+)^n\), and this lies in \(P\).  Since \(n\) is
odd, \((k_r^+)^n\equiv k_r^+\pmod {4O}\).  The detector in \eqref{eq:163-trace-detector}
factors through this layer and kills exactly \(QP\).
This proves \eqref{eq:163-mixed-ratio-detector}.
\end{proof}

\paragraph{Scope.}
This proposition uses the full-GT beta cocycle, its coefficient
equivariance and unitarity, and the degree--163 scalar and
first-layer identifications.  It does not use an arbitrary choice of
an adelic gamma function or an extension to motivic Galois groups.
The canonical factor is constructed directly on the finite cyclic
group.  The equality says that a particular off-diagonal ratio does
not lose the unknown residual; it supplies no new reason for that
residual to vanish.

\subsection{A three-point M\"obius orbit of the square roots}

Retain the periods \(\theta_i\) of
Lemma~\ref{lem:163-period-unit-square}, numbered by multiplication
of their cosets by two, and put
\[
 v_0=5+\theta_1,\qquad v_1=5+\theta_2,\qquad
 v_2=5+\theta_0.
\]
The period identity is \(U_i=-v_i^2\).

\begin{lemma}[An explicit special fiber]
\label{lem:163-mobius-period-fiber}
The three units satisfy
\[
 v_iv_{i+1}=v_i-1,\qquad
 v_{i+1}=1-v_i^{-1}\quad(i\bmod3),
\]
and are the roots of
\[
 X^3-14X^2+11X+1.
\]
In particular they lie in the fiber above \(14\) of the
degree-three quotient map
\[
 J(v)=\frac{v^3-3v+1}{v^2-v}
\]
for the order-three transformation \(v\mapsto1-v^{-1}\).
The discriminant of the polynomial is \(163^2\).
\end{lemma}

\begin{proof}
Apply \(s\) to \eqref{eq:163-period-table} to get
\[
 \theta_1\theta_2=17\theta_1+16\theta_2+21\theta_0.
\]
Using \(\theta_0+\theta_1+\theta_2=-1\), one obtains
\[
 (5+\theta_1)(5+\theta_2)
 =25+22\theta_1+21\theta_2+21\theta_0
 =4+\theta_1=v_0-1.
\]
Its two conjugates give the other equalities.  Division is allowed
because all three \(v_i\) are units.  Their sum is \(14\).
Summing the pair products gives \(11\), and multiplying the
Mobius relations, or using \(U_0U_1U_2=-1\), gives their product
\(-1\).  More directly the first two Mobius relations give
\(v_2=-1/(v_0-1)\), so \(v_0v_1v_2=-1\).
These elementary symmetric functions give the cubic polynomial.

The transformation \(\varphi(v)=1-v^{-1}\) has order three, and
\[
 v+\varphi(v)+\varphi^2(v)
   =v+\frac{v-1}{v}-\frac1{v-1}
   =\frac{v^3-3v+1}{v^2-v}.
\]
Thus \(J\) is invariant under \(\varphi\), has degree three,
and is its quotient map.  Setting \(J(v)=14\) gives exactly
the cubic above.  The elementary cubic discriminant formula gives
\(26569=163^2\).
\end{proof}

\paragraph{What additional compatibility would be needed.}
The relations express the square roots by the standard order-three
permutation of the three punctures of the thrice-punctured line.
This is a concrete geometric structure beyond an unspecified real
unit.  Applying it requires preservation of the particular fiber
\(J=14\).  An arithmetic identification
of that fiber with the cubic period field does not itself identify
the two actions of full GT on its relevant finite covers.

A sufficient next input would be a proved full-GT comparison between
the Kummer covers of this \(J=14\) fiber and the cyclotomic
\(163\)-cover defining the three \(U_i\), including the chosen
base paths and their signs.  Such a comparison would permit the
relations \(U_i=-v_i^2\) to be applied to the full-GT first-layer
coordinates.  The two- and three-cycle equations supply the Mobius
symmetry, but have not here been shown to supply preservation of
this particular fiber or its arithmetic identification.  No such comparison has been established here.
